\chapter{Compiler Architecture}
\label{chap:architecture}

This chapter describes the architecture of the \epl{} compiler as implemented in
the repository. The compiler is a Rust workspace with two crates. The main crate,
\texttt{epl}, contains the compiler and command-line interface. The helper crate,
\texttt{epl\_test}, contains the snapshot and smoke-test harness.

\section{Repository Structure}

The source tree is intentionally small. The main compiler modules are located in
\texttt{src/}. Documentation is located in \texttt{doc/}. Example programs and
their expected \irtree{} snapshots are located in \texttt{examples/}. Additional
regression tests are located in \texttt{tests/}.
Table~\ref{tab:repo-files} lists the main implementation files and their roles.

\begin{table}[H]
  \centering
  \footnotesize
  \begin{tabular}{p{0.40\textwidth}p{0.45\textwidth}}
    \toprule
    Path & Purpose \\
    \midrule
    \texttt{src/lex.rs} & Lexer, tokens, spans, and lexical errors. \\
    \texttt{src/ast.rs} & AST definitions and recursive descent parser. \\
    \texttt{src/ir\_tree.rs} & Typed tree IR data structures and module construction. \\
    \texttt{src/ir\_tree/lower\_ast.rs} & Semantic lowering from AST to \irtree{}. \\
    \texttt{src/ir\_tree/checkers.rs} & Main signature, purity, and compile-time checks. \\
    \texttt{src/ir\_tree/evaluator.rs} & Compile-time evaluator. \\
    \texttt{src/ir.rs} & Lower IR data structures. \\
    \texttt{src/ir/lower\_ir\_tree.rs} & Lowering from \irtree{} to CFG IR. \\
    \texttt{src/ir/opt/} & Lower IR cleanup passes. \\
    \texttt{src/llvm.rs} & LLVM module construction and object emission. \\
    \texttt{src/diagnostics.rs} & Shared diagnostic printing. \\
    \texttt{epl\_test/src/main.rs} & Test harness over examples and tests. \\
    \bottomrule
  \end{tabular}
  \caption{Main implementation files in the \epl{} repository}
  \label{tab:repo-files}
\end{table}

\section{Command-Line Interface}

The compiler executable accepts a command and a source file. The command chooses
which pipeline stage is printed or emitted. This is useful both for users and for
compiler development because each representation can be inspected independently.

\begin{lstlisting}[caption={Compiler command-line modes}]
epl ast <file>
epl ir_tree <file>
epl ir <file>
epl cfg <file>
epl llvm-ir <file>
epl llvm-obj <file>
\end{lstlisting}

\texttt{ast} prints the parsed AST. \texttt{ir\_tree} prints the typed tree.
\texttt{ir} prints the lower IR debug representation. \texttt{cfg} prints a
Graphviz DOT graph of the lower IR control-flow graph. \texttt{llvm-ir} prints
the generated LLVM module. \texttt{llvm-obj} asks LLVM to emit a native object
file named \texttt{a.out.o}.

The commands are not separate compilers. They call the same staged functions:
parse to AST, lower to \irtree{}, lower to IR, build an LLVM module, and
optionally emit an object file. If an earlier stage fails, diagnostics are
printed and the process exits.

\section{Pipeline Overview}

The full compiler pipeline can be summarized as:

\begin{lstlisting}[caption={Compiler pipeline}]
source text
  -> lexical tokens
  -> AST
  -> typed IR_TREE
  -> evaluated comptime constants
  -> lower CFG IR
  -> cleaned CFG IR
  -> LLVM IR
  -> native object file
\end{lstlisting}

Each stage has a deliberately narrow responsibility. The lexer does not know
about types. The parser does not resolve variable names. \irtree{} lowering
performs most semantic checks. The compile-time evaluator operates on typed
\irtree{} expressions rather than the AST. Lower IR generation makes control
flow and memory operations explicit. LLVM generation maps this lower IR to the
LLVM C API.

This separation is important for debugging. When an example does not compile,
the developer can inspect the AST to see whether parsing was correct, inspect
the \irtree{} snapshot to see whether semantic lowering was correct, inspect the
lower IR to understand CFG construction, or inspect LLVM IR to debug backend
generation.

\section{Entity Identifiers}

Source names are convenient for programmers but inconvenient for compiler
internals. A local variable can shadow another local variable with the same
name. A loop may contain nested loops. A function declaration may be referenced
before its body appears. The compiler therefore uses generated entity IDs for
functions, variables, and loops after semantic lowering.

The \texttt{make\_entity\_id!} macro creates small strongly typed wrappers around
non-zero integer IDs. The typed wrappers prevent accidentally mixing a function
ID with a variable ID or loop ID. They also provide compact debug formatting,
such as \texttt{fn\_1}, \texttt{var\_3}, and \texttt{loop\_2}. This is a small
implementation detail, but it improves the readability of generated snapshots.

\section{Type System Context}

The type system stores built-in types, struct definitions, and interned type
IDs. Struct and array types need layout information. The compiler computes each
struct's field offsets, size, and alignment when the struct definition is
processed. Array layout is computed from the element layout and length. Pointer
layout uses the target pointer size stored in the type system context. The
current implementation uses an 8-byte pointer size, matching common 64-bit
targets.

The type system uses a distinction between a \texttt{Type} value and a
\texttt{TypeId}. Many types can be stored directly in a compact enum, but typed
pointers and arrays need stable references to pointee or element types.
\texttt{TypeId} provides this stable reference. For example, \texttt{*T} stores
the ID of \texttt{T}, and \texttt{[T; N]} stores the ID of \texttt{T} plus the
array length.

\section{Diagnostics}

The compiler stages expose different error types, but they all adapt to a shared
diagnostic interface. The diagnostic printer displays the message, file path,
line number, source line, and a caret underline for the span. This is why the
lexer and parser retain byte offsets and why later semantic errors carry spans
from AST nodes.

Good diagnostics matter even in a small compiler. A type error without a source
location is difficult to fix. A parser error that does not say what token was
expected is frustrating. The current diagnostic system is simple, but it already
captures the most important user-facing behavior: an error points to the source
fragment that caused it.

\section{Documentation and Examples}

The repository includes four documentation files: syntax reference, language
reference, compiler architecture, and LLVM setup. These files serve two
purposes. First, they document the user-visible language rules. Second, they
provide a specification against which the implementation and thesis can be
checked.

The examples cover hello world, recursion, iterative loops, arrays, structs,
pointers, external \texttt{exit}, \texttt{else if}, and compile-time prime
generation. Each example has a neighboring \irtree{} snapshot. The snapshots
are not only tests. They are executable documentation of lowering decisions,
such as how \texttt{for} loops become hidden counter variables or how a
\texttt{comptime} block becomes a constant array.

\section{Build Environment}

The Rust workspace uses edition 2024 and depends on \texttt{llvm-sys} version
\texttt{221.0.0}. The project includes a Cargo configuration file that points
\texttt{LLVM\_SYS\_221\_PREFIX} to an LLVM 22.1.2 installation. The compiler can
be built with:

\begin{lstlisting}[caption={Building the compiler}]
cargo build
\end{lstlisting}

The test harness can then be run with:

\begin{lstlisting}[caption={Running the EPL test harness}]
cargo run -p epl_test ./target/debug/epl examples tests
\end{lstlisting}

The command in Listing above is also used in Chapter \ref{chap:testing} to
report the verification status of the implementation.
